\chapter{Mitral valve surgery}
\index{Mitral valve surgery}

\section*{Definition and Overview}
Mitral valve surgery comprises a range of surgical procedures performed to treat dysfunction of the mitral valve, which is situated between the left atrium and left ventricle of the heart. The primary disorders necessitating surgery are mitral regurgitation (where the valve does not close tightly, causing backward blood flow) and mitral stenosis (where the valve opening is narrowed, restricting forward blood flow). Surgical interventions are categorised into mitral valve repair (preserving the patient's native valve tissue using techniques like annuloplasty, chordal reconstruction, or leaflet resection) and mitral valve replacement (replacing the diseased valve with a bioprosthetic tissue valve or a mechanical valve). Mitral valve repair is clinically preferred due to its superior long-term survival, lower risk of thromboembolism and endocarditis, and avoidance of lifelong anticoagulation therapy.

\section*{Epidemiology}
Valvular heart disease is a significant contributor to cardiovascular morbidity. In many countries and other developed countries, degenerative mitral valve disease (such as myxomatous degeneration or mitral valve prolapse) is the most common indication for surgery, with a rising prevalence corresponding to an ageing population. Mitral regurgitation affects approximately 2\% of the general population, rising to over 10\% in individuals over 75 years of age. Conversely, mitral stenosis is primarily caused by rheumatic heart disease (RHD), which remains relatively rare in the general the population but is highly prevalent within Aboriginal and Torres Strait Islander communities and immigrants from developing nations. Mitral valve surgery represents one of the most frequently performed valvular procedures globally, secondary only to aortic valve replacement.

\section*{Aetiology and Pathophysiology}
Mitral valve dysfunction requiring surgical intervention develops from primary structural abnormalities or secondary ventricular changes.
\begin{itemize}
    \item \textbf{Degenerative disease:} Myxomatous degeneration leading to mitral valve prolapse, leaflet flail, or chordal rupture, preventing proper leaflet coaptation.
    \item \textbf{Rheumatic heart disease:} Chronic inflammation and subsequent scarring, fusion, and calcification of the valve leaflets and chordae tendineae, predominantly resulting in stenosis.
    \item \textbf{Ischaemic or functional regurgitation:} Ventricular remodelling, dilation, or papillary muscle displacement caused by coronary artery disease or dilated cardiomyopathy, pulling the valve leaflets apart.
    \item \textbf{Infective endocarditis:} Bacterial or fungal infection destroying valve leaflets, chordae, or the annulus, leading to acute, severe regurgitation.
    \item \textbf{Congenital malformations:} Rarest aetiology, presenting as a parachute mitral valve or cleft mitral valve in paediatric populations.
\end{itemize}

\section*{Clinical Features}
The clinical presentation of mitral valve disease depends on whether the condition is acute (e.g., chordal rupture or endocarditis) or chronic and progressive.
\begin{itemize}
    \item \textbf{Exertional dyspnoea:} Shortness of breath during physical activity, progressing to orthopnoea (breathlessness when lying flat) and paroxysmal nocturnal dyspnoea.
    \item \textbf{Fatigue and exercise intolerance:} Reduced cardiac output leading to persistent tiredness and inability to complete everyday tasks.
    \item \textbf{Palpitations:} Sensation of a racing or irregular heartbeat, frequently associated with the development of atrial fibrillation due to left atrial enlargement.
    \item \textbf{Peripheral oedema:} Swelling of the ankles, feet, or abdomen, indicating secondary right-sided heart failure.
    \item \textbf{Auscultatory findings:} A characteristic pansystolic murmur radiating to the axilla in mitral regurgitation, or a mid-diastolic rumble with an opening snap in mitral stenosis.
\end{itemize}

\section*{Differential Diagnosis}
Symptoms of mitral valve disease can overlap significantly with other cardiovascular and pulmonary pathologies.
\begin{itemize}
    \item \textbf{Aortic valve disease:} Presents with similar exertional dyspnoea and fatigue, requiring echocardiographic differentiation.
    \item \textbf{Coronary artery disease:} Can cause dyspnoea and chest pain; may co-exist with or directly cause functional mitral regurgitation.
    \item \textbf{Congestive heart failure:} Left ventricular failure leading to dyspnoea and pulmonary congestion without primary valvular disease.
    \item \textbf{Chronic obstructive pulmonary disease (COPD):} Causes progressive dyspnoea and exercise limitation, but is distinguished by pulmonary function testing.
    \item \textbf{Pulmonary hypertension:} Can be primary, though it is frequently secondary to chronic mitral valve disease.
\end{itemize}

\section*{When to Seek Medical Care}
Patients with known mitral valve disease should undergo regular clinical and echocardiographic surveillance. They should seek medical attention if they experience worsening shortness of breath, sudden palpitations, or ankle swelling.

\textbf{Contact your local emergency services for:}
\begin{itemize}
    \item Severe, sudden-onset shortness of breath (suggestive of acute pulmonary oedema).
    \item Crushing chest pain or pressure radiating to the neck, jaw, or left arm.
    \item Sudden loss of consciousness, fainting, or severe dizziness.
    \item Symptoms of a stroke, such as sudden weakness or numbness on one side of the body or difficulty speaking.
\end{itemize}

\section*{Diagnosis and Investigations}
A comprehensive diagnostic workup is essential to determine the severity of valve disease, assess ventricular function, and plan the surgical approach.
\begin{itemize}
    \item \textbf{Transthoracic echocardiography (TTE):} The primary non-invasive imaging modality to assess valve anatomy, quantify stenosis or regurgitation, and evaluate left ventricular ejection fraction.
    \item \textbf{Transoesophageal echocardiography (TOE):} Provides high-resolution, detailed images of the valve leaflets and subvalvular apparatus, crucial for assessing the feasibility of repair immediately before or during surgery.
    \item \textbf{Electrocardiogram (ECG):} Detects signs of left atrial enlargement, left ventricular hypertrophy, and arrhythmias, particularly atrial fibrillation.
    \item \textbf{Coronating angiography:} Performed preoperatively in patients over 40 or with risk factors to identify concurrent coronary artery disease requiring bypass grafting.
    \item \textbf{Cardiac MRI:} Indicated when echocardiographic assessment is inconclusive regarding the severity of regurgitation or ventricular volumes.
    \item \textbf{Chest X-ray:} May show cardiomegaly, left atrial enlargement, and signs of pulmonary venous congestion.
\end{itemize}

\section*{Management}
Surgical management is indicated for symptomatic patients or asymptomatic patients with evidence of left ventricular dysfunction, pulmonary hypertension, or new-onset atrial fibrillation.
\begin{itemize}
    \item \textbf{Mitral valve repair:} The gold standard surgical treatment, employing annuloplasty rings to stabilise the valve annulus, along with leaflet resection or artificial chordae replacement.
    \item \textbf{Mitral valve replacement:} Reserved for cases where repair is technically impossible; involves replacing the valve with a mechanical prosthesis (requiring lifelong warfarin therapy) or a bioprosthetic valve (which degenerates over 10--15 years but avoids long-term anticoagulation).
    \item \textbf{Minimally invasive surgery:} Accessing the heart via a mini-thoracotomy or robotic assistance, reducing recovery times and surgical trauma compared to a full median sternotomy.
    \item \textbf{Transcatheter edge-to-edge repair (TEER):} A percutaneous, catheter-based approach (e.g., MitraClip) for high-surgical-risk or inoperable patients with severe mitral regurgitation.
    \item \textbf{Medical optimisation:} Pre- and post-operative pharmacotherapy using diuretics (to manage fluid overload), beta-blockers (for rate control), and anticoagulation.
\end{itemize}

\section*{Complications}
Mitral valve surgery is a major cardiothoracic procedure and carries risks of immediate post-operative and long-term complications.
\begin{itemize}
    \item \textbf{Post-operative haemorrhage:} Bleeding requiring re-exploration of the chest or blood transfusions.
    \item \textbf{Cardiovascular events:} Stroke, myocardial infarction, transient ischaemic attacks, or low cardiac output syndrome.
    \item \textbf{Arrhythmias:} High incidence of post-operative atrial fibrillation, as well as heart block requiring a permanent pacemaker.
    \item \textbf{Thromboembolism and anticoagulation complications:} Valve thrombosis or systemic embolisation, contrasted with haemorrhagic risks of warfarin therapy (primarily associated with mechanical valves).
    \item \textbf{Prosthetic valve dysfunction:} Endocarditis (infection of the prosthetic valve) or structural valve degeneration over time (in bioprosthetic valves).
\end{itemize}

\section*{Prognosis and Outcomes}
The prognosis following mitral valve surgery is generally excellent, particularly when mitral valve repair is performed before the onset of irreversible left ventricular dysfunction or severe pulmonary hypertension. The 10-year survival rate following successful mitral valve repair exceeds 80--90\%, matching the age-adjusted survival of the general population. In contrast, replacement carries slightly lower long-term survival rates and higher complication rates due to prosthesis-related issues. The recovery period after open-heart surgery typically ranges from 6 to 12 weeks, with significant improvements in exercise capacity, symptoms, and overall quality of life.

\section*{Prevention and Self-Care}
Post-operative recovery and long-term care require active patient participation and adherence to medical guidelines.
\begin{itemize}
    \item \textbf{Anticoagulation compliance:} Strict adherence to daily warfarin therapy, including regular blood tests (INR monitoring), for patients with mechanical valves.
    \item \textbf{Endocarditis prophylaxis:} Taking preventive antibiotics before certain dental or invasive procedures, as recommended by cardiologists, to prevent prosthetic valve infection.
    \item \textbf{Cardiac rehabilitation:} Participating in a structured cardiac rehabilitation programme to safely restore physical strength and cardiovascular endurance.
    \item \textbf{Lifestyle modifications:} Adhering to a heart-healthy diet, managing blood pressure, stopping smoking, and maintaining a healthy weight.
    \item \textbf{Infection monitoring:} Frequently inspecting the surgical wound for redness, swelling, or drainage, and monitoring for fever.
\end{itemize}

\section*{Sources and Further Reading}
The following reputable organisations provide further information and resources regarding mitral valve surgery:
\begin{itemize}
    \item American Heart Association. \textit{Mitral Valve Repair or Replacement.} https://www.heart.org/
    \item Cleveland Clinic. \textit{Mitral Valve Surgery.} https://my.clevelandclinic.org/
    \item National Heart Foundation of Australia. \textit{Heart valve surgery.} https://www.heartfoundation.org.au/
    \item NHS. \textit{Mitral valve surgery.} https://www.nhs.uk/
    \item Society of Thoracic Surgeons. \textit{Mitral Valve Repair and Replacement.} https://www.sts.org/
\end{itemize}
